\documentclass{article}

\usepackage{amsmath,amssymb}
\usepackage{xurl}
\usepackage[hidelinks]{hyperref}
\setlength{\emergencystretch}{2em}

% Shared commands for notes in docs/paper-gaps/.

\DeclareUnicodeCharacter{2080}{\ifmmode{}_{0}\else\textsubscript{0}\fi}
\DeclareUnicodeCharacter{2081}{\ifmmode{}_{1}\else\textsubscript{1}\fi}
\DeclareUnicodeCharacter{2082}{\ifmmode{}_{2}\else\textsubscript{2}\fi}
\DeclareUnicodeCharacter{2099}{\ifmmode{}_{n}\else\textsubscript{n}\fi}

\newcommand{\Lean}{\textsc{Lean}}
\ExplSyntaxOn
\NewDocumentCommand{\leanid}{m}{
  \ifmmode
    \textsf{\detokenize{#1}}
  \else
    \group_begin:
    \urlstyle{sf}
    \tl_set:Nn \l_tmpa_tl {#1}
    \tl_replace_all:Nnn \l_tmpa_tl {\_} {_}
    \exp_args:NV \path \l_tmpa_tl
    \group_end:
  \fi
}
\ExplSyntaxOff
\providecommand{\tr}{\operatorname{tr}}
\newcommand{\tnleanIssueBase}{https://github.com/LionSR/TNLean/issues/}
\newcommand{\tnleanPrBase}{https://github.com/LionSR/TNLean/pull/}
\newcommand{\tnleanDocBase}{https://sirui-lu.com/TNLean/docs/}
\newcommand{\ghissue}[1]{\href{\tnleanIssueBase#1}{GitHub issue \##1}}
\newcommand{\ghpr}[1]{\href{\tnleanPrBase#1}{PR \##1}}
\newcommand{\doclink}[1]{\href{\tnleanDocBase#1.html}{\path{#1}}}

% Dirac notation and the identity operator, as in blueprint/src/macros/common.tex.
% \providecommand keeps note-local definitions authoritative.
\usepackage{dsfont}
\providecommand{\ket}[1]{|#1\rangle}
\providecommand{\bra}[1]{\langle #1|}
\providecommand{\braket}[2]{\langle #1 | #2 \rangle}
\providecommand{\ketbra}[2]{|#1\rangle\!\langle #2|}
\providecommand{\Id}{\mathds{1}}
\providecommand{\one}{\mathds{1}}
\providecommand{\C}{\mathbb{C}}
\providecommand{\MN}[1]{M_{#1}(\C)}
\providecommand{\MD}{M_D(\C)}

% Verdict marker: \gapnote{<kind>}{<status>}. Kind is the policy
% classification (clarification, local-correction, scope-restriction,
% unfaithful, false-source, open-gap); status is open, wip (elimination
% underway), resolved, or historical. Machine-read by the site index;
% prints a verdict line.
\newcommand{\gapnote}[2]{\par\noindent\textbf{Verdict:} #1 (#2).\par}


\title{The norm of the frustration-free three-region replacement in ALVV}
\author{}
\date{November 4, 2026}
\begin{document}
\maketitle
\gapnote{local-correction}{resolved}

\paragraph{Source.}
The proof of \cite[Theorem~5.15]{Arad2016RigorousRG} (Theorem \texttt{t:aisff},
lines 1592--1607) replaces the three bulk regions to the left of, between and to
the right of two windows of width about $2\ell$ by the regional construction of
\cite[Definition~5.10]{Arad2016RigorousRG}. It then states, at line 1599, that the
replaced Hamiltonian $\widetilde H$ ``has norm $O(1)$'', and, at line 1606, that
each improving operator $A_i$ has bond dimension $O(d^k)$, where $k$ is the
degree of the Chebyshev filter.

\paragraph{The norm statement fails.}
The bonds in the two windows are not replaced, so $\widetilde H$ contains about
$4\ell$ untruncated terms. Take qubits and
$h_i=I-|00\rangle\langle00|$ on the qubits $i,i+1$. These are commuting
orthogonal projections, diagonal in the computational basis. The energy of a
basis configuration is the number of bonds not in the state $00$. Hence
$|0\cdots0\rangle$ is the unique ground state, with energy $0$, and every other
configuration has energy at least $1$. The chain is frustration free, with a
unique ground state and gap $1$.
On $|1\cdots1\rangle$ each window bond acts as $1$. For every regional
parity layer, the product of the factors $I-h_i=|00\rangle\langle00|$ over a
nonempty set of bonds annihilates $|1\cdots1\rangle$, so the layer term
$I-P$ acts as $1$. With two windows of $2\ell$ bonds each and three regions
containing bonds of both parities,
\begin{align}
    \widetilde H\,|1\cdots1\rangle & =(4\ell+6)\,|1\cdots1\rangle.
    \label{eq:alvv_ff_norm_example}
\end{align}
Thus $\|\widetilde H\|=4\ell+6$, which grows with $\ell$.

\paragraph{The conclusion is unaffected.}
The parameter choice $\ell=\widetilde\Theta(\gamma^{-1}\log^2d)$ of the source
is the one that works with $\|\widetilde H\|=\Theta(\ell)$. With the Chebyshev
filter of degree $k=\ell^2$, the gap $\eta_1=\Theta(\gamma)$ and the upper bound
$4\ell+6$, the filter contracts the excited spectrum by
$\Delta=4\exp(-\Omega(\ell^{3/2}\gamma^{1/2}))$, while the number of improving
operators is $D^2$ with $\log D=O(\ell\log(d\ell))$. The trade-off
$D^{12}\Delta<10^{-5}$ therefore requires $\ell\gtrsim\gamma^{-1}\log^2(d\ell)$,
the printed choice. With norm $O(1)$ the smaller choice
$\ell=\widetilde\Theta(\gamma^{-1/2}\log d)$ would have sufficed, and the
source would have obtained a better bound on $D$. The formalization uses the
bound $\|\widetilde H\|\le4\ell+6$, proved in
\leanid{TNLean.MPS.AreaLaw.ALVV.norm_threeRegion_le}.

\paragraph{Bond dimension.}
The replaced Hamiltonian has operator Schmidt rank of order $d^2$ at every cut,
not of order $d$. The argument of the source therefore gives bond dimension of
order $(d^2+O(1))^{k+O(1)}$ for the improving operators, not $O(d^k)$. The
formalization proves the bound $(d^2+4)^{\ell^2+1}$ at every cut inside the
prescribed region, with logarithm $O(u^5/g^2)$, $u=1+\log d+\log(1/g)$,
$g=\min(\gamma,1)$. This is the form $2^{\widetilde O(\gamma^{-2}\log^5d)}$ of
the theorem statement, which is unaffected.

\paragraph{Formalized statement.}
\leanid{TNLean.MPS.AreaLaw.ALVV.exists_frustrationFree_improving_family}
proves the existence part of Theorem~5.15 with these bounds. It needs neither
uniqueness of the ground state nor frustration freeness, and it does not assume
that the terms are projections: each term is replaced by its range projection,
as at line 1457 of the source. The improvement is stated for viability error
$\delta<1$, the boundary convention already used for Theorem~4.7. The running
time is not formalized.

\bibliographystyle{alpha}
\bibliography{references}

\end{document}
